\documentclass[11pt]{article}

\usepackage[a4paper,margin=1in]{geometry}
\usepackage{amsmath,amssymb,amsthm,mathtools}
\usepackage{booktabs}
\usepackage{array}
\usepackage{hyperref}
\usepackage{xcolor}

\hypersetup{
  colorlinks=true,
  linkcolor=blue!55!black,
  citecolor=blue!55!black,
  urlcolor=blue!55!black
}
\setlength{\emergencystretch}{3em}

\newtheorem{theorem}{Theorem}
\newtheorem{lemma}{Lemma}
\newtheorem{proposition}{Proposition}
\newtheorem{corollary}{Corollary}
\newtheorem{remark}{Remark}

\newcommand{\C}{\mathbb{C}}
\newcommand{\Spec}{\operatorname{Spec}}
\newcommand{\QLS}{\operatorname{QLS}}
\newcommand{\supp}{\operatorname{supp}}
\newcommand{\Span}{\operatorname{span}}
\newcolumntype{L}[1]{>{\raggedright\arraybackslash}p{#1}}

\title{Low Cardinalities of Quantum Latin Squares of Order Five:\\
Computational Certificates for the Gaps at 8, 9, 10, and 11}
\author{Research notes generated from the QLS spectrum project}
\date{June 12, 2026}

\begin{document}
\maketitle

\begin{abstract}
We record the current small-order evidence for the cardinality spectrum of
quantum Latin squares.  The main new computational conclusion is that no
quantum Latin square of order five has cardinality \(8\), \(9\), \(10\), or
\(11\).  Together
with the general exclusion of cardinality \(n+1\), and with explicit
constructions of cardinalities \(5,7,12\), this gives
\[
  \{5,7,12\}\subseteq \Spec(5),\qquad
  6,8,9,10,11\notin \Spec(5).
\]
The exclusions of \(8\), \(9\), \(10\), and \(11\) are finite, reproducible,
computer-assisted classifications.  The final obstruction checks are ordinary
two- and three-dimensional linear algebra.  We also record an exact
cardinality-\(12\) construction and the next target \(c=13\).
\end{abstract}

\section{Definitions and context}

A quantum Latin square of order \(n\), abbreviated \(\QLS(n)\), is an
\(n\times n\) array \(Q=(q_{ij})\) whose entries are unit vectors in
\(\C^n\), such that each row and each column is an orthonormal basis of
\(\C^n\).  Two entries are identified if they differ by a nonzero scalar; in
the unit-vector normalization this means they differ by a global phase.  The
cardinality of \(Q\) is the number of distinct projective classes appearing in
the array.

Let \(\Spec(n)\) denote the set of possible cardinalities of \(\QLS(n)\).
The trivial bounds are
\[
  n \le c \le n^2.
\]
Classical Latin squares give \(c=n\).  General constructions give
maximal cardinality \(n^2\) for all \(n\ge 4\).  Zhang--Wang--Ji proved that
for all multiples \(n=4m\), \(m\ge 2\),
\[
  \Spec(n)=[n,n^2]\setminus\{n+1\}.
\]
They also proved the following general obstruction.

\begin{lemma}[Zhang--Wang--Ji]
For every \(n\ge 2\), no \(\QLS(n)\) has cardinality \(n+1\).
\end{lemma}

Thus \(6\notin\Spec(5)\).  The remaining low-cardinality question for order
five begins at \(c=7,8,9,\ldots\).

\section{Main order-five conclusions}

\begin{theorem}[Certified low spectrum of order five]
The current certified low-cardinality data for order five is
\[
  \{5,7,12\}\subseteq \Spec(5),\qquad
  6,8,9,10,11\notin \Spec(5).
\]
\end{theorem}

\begin{proof}
The value \(5\) is classical.  The value \(7\) is obtained by rotating a
single intercalate in a classical Latin square of order five.  The value
\(12\) is given explicitly in Section~\ref{sec:c12}.  The exclusion of \(6\)
is the \(n+1\) lemma.  The exclusions of \(8\) and \(9\) are the finite
certificates described in Sections~\ref{sec:c8} and~\ref{sec:c9}; the
exclusion of \(10\) is described in Section~\ref{sec:c10}; the exclusion of
\(11\) is described in Section~\ref{sec:c11}.
\end{proof}

\begin{corollary}
The naive extension
\[
  \Spec(n)=[n,n^2]\setminus\{n+1\}\quad\text{for all }n\ge 5
\]
is false already at \(n=5\).
\end{corollary}

\section{Normalized pattern method}

Let \(Q\) be a hypothetical \(\QLS(5)\) of small cardinality.  Applying a
global unitary, we may assume that the first row is
\[
  e_0,e_1,e_2,e_3,e_4.
\]
The old projective classes are labelled \(0,1,2,3,4\), corresponding to the
coordinate vectors.  If the cardinality is \(5+t\), the remaining classes are
labelled \(5,\ldots,4+t\).  A cell below the first row is called active if it
contains a new label.

No label can repeat in a row or column.  In addition, old label \(j\) cannot
occur below the first-row entry in column \(j\).  For a new label \(k\), define
\[
  Z_k=\{j\in\{0,1,2,3,4\}: x_k\perp e_j \text{ is forced}\},\qquad
  S_k=\{0,1,2,3,4\}\setminus Z_k.
\]
The vector \(x_k\) must lie in \(\Span\{e_j:j\in S_k\}\).  Since \(x_k\) is
not a coordinate vector, one must have
\[
  |S_k|\ge 2.
\]
Whenever two new labels occur in a common row or column, their vectors must be
orthogonal.  This gives an orthogonality graph on the new labels.

\section{The gap at cardinality 8}
\label{sec:c8}

\begin{theorem}
There is no \(\QLS(5)\) of cardinality \(8\).
\end{theorem}

\begin{proof}
After the first-row normalization there are three new labels, \(5,6,7\).
The deterministic classifier enumerates all admissible normalized label
patterns and records the support sets \(S_5,S_6,S_7\) and the forced
orthogonality graph.  It gives the following exhaustive classification.

\begin{center}
\begin{tabular}{@{}L{0.33\linewidth}rL{0.43\linewidth}@{}}
\toprule
Obstruction type & Patterns & Explanation \\
\midrule
\texttt{same\_2d\_path} & 8640 &
three vectors in one \(2\)-space with a path \\
\texttt{same\_2d\_triangle} & 23040 &
three mutually orthogonal vectors in one \(2\)-space \\
\texttt{two\_2d\_one\_3d\_triangle} & 2880 &
two vectors form a \(2\)-space basis; the third degenerates \\
\bottomrule
\end{tabular}
\end{center}

The full certificate statistics are
\[
\begin{array}{rl}
\text{unique patterns} &= 34560,\\
\text{unobstructed patterns} &= 0,\\
\text{SHA-256 digest} &=
\begin{gathered}
\texttt{8837bc4dd3c92ab424d30c4f1be7f20f}\\
\texttt{235de8e87cee6563d6c202a609ed1a75}
\end{gathered}.
\end{array}
\]
The three obstruction types are impossible by elementary linear algebra.
First, a two-dimensional Hilbert space cannot contain three mutually
orthogonal nonzero projective vectors.  Second, if \(x_5\perp x_6\) and
\(x_6\perp x_7\) inside the same two-dimensional subspace, then \(x_5\) and
\(x_7\) lie in the same one-dimensional orthogonal complement of \(x_6\), so
they are projectively equal.  Third, if two vectors form an orthogonal basis
of a two-dimensional support \(S\), then any vector orthogonal to both has
zero projection onto \(S\); if its support is \(S\cup\{t\}\), it is forced onto
the coordinate line \(\C e_t\), so it is not a new projective class.
\end{proof}

\section{The gap at cardinality 9}
\label{sec:c9}

\begin{theorem}
There is no \(\QLS(5)\) of cardinality \(9\).
\end{theorem}

\begin{proof}
There are four new labels, \(5,6,7,8\).  The active-cell graph has four row
vertices, five column vertices, and an edge for each active cell.  The script
\texttt{qls5\_active\_graph\_general.py --t 4} enumerates all \(2^{20}\)
active graphs up to row and column permutations and applies the basic
line-degree, old-label fillability, new-label colorability, and local support
conditions.  The counts are:
\[
\begin{array}{rl}
\text{raw graphs} &= 1048576,\\
\text{line-degree feasible types} &= 151,\\
\text{old-fill feasible types} &= 151,\\
\text{new-coloring feasible types} &= 150,\\
\text{old-and-new feasible types} &= 150,\\
\text{local-support feasible types} &= 14.
\end{array}
\]
Thus only \(14\) active graph types survive.

The deterministic certificate then expands these \(14\) graph types into all
normalized old/new label fillings.  It finds
\[
\begin{array}{rl}
\text{label patterns} &= 1587,\\
\text{support-edge types} &= 97,\\
\text{unclassified support-edge types} &= 0.
\end{array}
\]
All \(97\) support-edge types fall into the following three deterministic
obstruction classes.

\begin{center}
\begin{tabular}{lrr}
\toprule
Obstruction type & Support-edge types & Label patterns \\
\midrule
\texttt{clique\_union\_dim} & 47 & 1451 \\
\texttt{same\_2d\_basis\_common\_neighbor} & 33 & 44 \\
\texttt{same\_2d\_two\_step} & 17 & 92 \\
\bottomrule
\end{tabular}
\end{center}

The first obstruction says that a clique of \(q\) mutually orthogonal nonzero
vectors cannot be supported inside a coordinate subspace of dimension less
than \(q\).  The second says that if two adjacent vectors have the same
two-dimensional support, they form a basis of that support; a common neighbor
must therefore be supported outside that two-space, and if fewer than two
outside coordinates remain, it is zero or a coordinate vector.  The third is
the two-step path obstruction already used for \(c=8\).  Since every remaining
support-edge type is obstructed, no cardinality-\(9\) square exists.
\end{proof}

\section{The gap at cardinality 10}
\label{sec:c10}

\begin{theorem}
There is no \(\QLS(5)\) of cardinality \(10\).
\end{theorem}

\begin{proof}
There are five new labels, \(5,6,7,8,9\).  The active graph sieve starts with
all \(2^{20}\) active-cell graphs and applies the same line-degree,
old-label fillability, new-label colorability, and local support constraints.
The counts are
\[
\begin{array}{rl}
\text{raw graphs} &= 1048576,\\
\text{line-degree feasible types} &= 222,\\
\text{old-fill feasible types} &= 222,\\
\text{new-coloring feasible types} &= 220,\\
\text{old-and-new feasible types} &= 220,\\
\text{local-support feasible types} &= 40.
\end{array}
\]
Thus only \(40\) active graph types survive the combinatorial and local
support filters.

The deterministic certificate expands those \(40\) graph types into all
normalized old/new label fillings and collapses them to support-edge types.
It finds
\[
\begin{array}{rl}
\text{label patterns} &= 9010,\\
\text{support-edge types} &= 1724,\\
\text{unclassified support-edge types} &= 0.
\end{array}
\]
The support-edge types are distributed as follows.
\begin{center}
\begin{tabular}{lrr}
\toprule
Obstruction type & Support-edge types & Label patterns \\
\midrule
\texttt{clique\_union\_dim} & 1064 & 8052 \\
\texttt{two\_2d\_overlap\_one\_edge} & 601 & 817 \\
\texttt{same\_2d\_two\_step} & 46 & 116 \\
\texttt{same\_2d\_basis\_common\_neighbor} & 13 & 25 \\
\bottomrule
\end{tabular}
\end{center}

The first, third, and fourth obstruction classes are those used for
cardinality \(9\).  The new class is
\texttt{two\_2d\_overlap\_one\_edge}: if two adjacent vectors have
two-dimensional supports \(S_u,S_v\) with \(|S_u\cap S_v|=1\), then their
inner product has only one possible common coordinate contribution.  Hence
orthogonality forces at least one vector to vanish on the common coordinate,
so that vector is supported on a single coordinate line and is not a new
projective class.  Since all \(1724\) support-edge types are obstructed, no
cardinality-\(10\) square exists.
\end{proof}

\section{The gap at cardinality 11}
\label{sec:c11}

\begin{theorem}
There is no \(\QLS(5)\) of cardinality \(11\).
\end{theorem}

\begin{proof}
There are six new labels after the first-row normalization.  The fast
active-graph layer gives \(220\) old-and-new feasible active graph types.
Instead of using the older generic local-support prefilter, the incremental
certificate enumerates label fillings directly over these \(220\) graph types,
enforcing \(|Z_k|\le 3\) during assignment and introducing new labels in
canonical first-occurrence order.

The complete enumeration gives
\[
\begin{array}{rl}
\text{label patterns} &= 201297,\\
\text{support-edge types} &= 46808.
\end{array}
\]
After reclassification with the final obstruction library, all support-edge
types are obstructed:
\[
\begin{array}{rl}
\text{unclassified support-edge types} &= 0.
\end{array}
\]
The final counts are
\begin{center}
\begin{tabular}{lrr}
\toprule
Obstruction type & Support-edge types & Label patterns \\
\midrule
\texttt{clique\_union\_dim} & 36853 & 184669 \\
\texttt{two\_2d\_overlap\_one\_edge} & 9804 & 16346 \\
\texttt{same\_2d\_two\_step} & 131 & 262 \\
\texttt{same\_2d\_basis\_common\_neighbor} & 18 & 18 \\
\texttt{same\_3d\_two\_independent\_common\_neighbors} & 2 & 2 \\
\bottomrule
\end{tabular}
\end{center}

The last obstruction is new at \(c=11\).  Suppose two distinct new vectors
\(u,w\) have the same three-dimensional support \(S\), and suppose they have
two common neighbors \(a,b\) whose two-dimensional supports are contained in
\(S\), meet in one coordinate, and together cover \(S\).  Then \(a\) and \(b\)
are necessarily linearly independent unless one degenerates to a coordinate
vector.  Therefore the orthogonal complement of \(\Span(a,b)\) in
\(\Span(S)\) is one-dimensional.  Both \(u\) and \(w\) lie in this line, so
they are projectively equal, contradicting that they are distinct new classes.
This eliminates the final two support-edge types; hence no cardinality-\(11\)
square exists.
\end{proof}

\section{An exact cardinality-\texorpdfstring{\(12\)}{12} construction}
\label{sec:c12}

Let \(e_i\) be the standard coordinate vectors of \(\C^5\).  Define the
following unnormalized vectors:
\[
\begin{aligned}
a&=(0,-1,0, 2, 1),&
b&=(0, 1,0, 1,-1),&
c&=(0, 1,0, 0, 1),\\
d&=(0,-1,0, 2,-5),&
h&=(0, 2,0, 1, 0),&
f&=(0, 1,0,-2, 0),\\
g&=(0, 1,0, 0,-1).
\end{aligned}
\]
After normalizing these seven vectors, the array
\[
\begin{array}{ccccc}
e_0 & e_1 & e_2 & e_3 & e_4\\
a   & e_0 & b   & c   & e_2\\
d   & e_2 & a   & e_0 & h\\
h   & e_4 & e_0 & e_2 & f\\
e_2 & e_3 & c   & g   & e_0
\end{array}
\]
is a \(\QLS(5)\) of cardinality \(12\).  Verification is exact: all row and
column orthogonality checks reduce to integer dot products before
normalization.  The corresponding verifier is the script
\path{exact_qls5_c12_construction.py}.

\begin{remark}
A tangent-space test around this construction was also performed.  With the
first row fixed, the base construction has projective tangent nullity \(4\).
For every tested single-cell split of a repeated class below the first row,
the projective tangent nullity remains \(4\).  Thus this construction appears
first-order rigid against the simplest local splitting route toward
cardinality \(13\).  This is not a nonexistence theorem for \(c=13\).
\end{remark}

\section{Next target}

The next unresolved low value is \(c=13\).  The beginning of the order-five
spectrum now has the certified jump
\[
  5,\ 7,\ 12,\ldots
\]
The known cardinality-\(12\) construction is first-order rigid against simple
single-cell splitting, so the next target likely requires either a more global
modification of the \(c=12\) pattern or a different construction altogether.

The first active-graph layer for \(c=13\) has \(216\) old-and-new feasible
graph types.  Incremental enumeration has currently completed graph types
\(1\) through \(194\), with several million-pattern graph types already
processed.  Sampled large graph types, including one with \(1{,}092{,}990\)
support-edge types, are fully covered by the current obstruction library.
However, the complete \(c=13\) classification is not yet finished: the
bottleneck is now data engineering, since some individual graph certificates
are multi-gigabyte-scale if every support-edge type is written to JSON.  The
next version should classify support-edge types online and save only
obstruction counts, samples of unclassified types, and reproducibility hashes.

\section{Reproducibility}

The main files produced in this stage are listed below.
\begin{center}
\begin{tabular}{@{}L{0.48\linewidth}L{0.43\linewidth}@{}}
\toprule
File & Role \\
\midrule
\path{c8_classification_certificate.py} &
deterministic cardinality-\(8\) classifier \\
\path{c8_classification_certificate.json} &
compact cardinality-\(8\) certificate \\
\path{qls5_active_graph_general.py} &
active-graph sieve for \(c=8,9,10\) \\
\path{qls5_c9_deterministic_certificate.py} &
proof-oriented cardinality-\(9\) certificate \\
\path{qls5_c9_deterministic_certificate.json} &
cardinality-\(9\) certificate output \\
\path{qls5_c9_support_obstruction.py} &
support-edge obstruction classifier \\
\path{qls_support_obstruction_general.py} &
general support-edge obstruction classifier \\
\path{qls5_c10_deterministic_probe.py} &
proof-oriented cardinality-\(10\) certificate \\
\path{qls5_c10_deterministic_probe_complete.json} &
cardinality-\(10\) certificate output \\
\path{qls5_incremental_probe.py} &
canonical incremental probe used for \(c=11\) \\
\path{qls5_c11_incremental_old_and_new_ok_complete_reclassified.json} &
cardinality-\(11\) certificate output \\
\path{qls5_c13_progress_note.md} &
current cardinality-\(13\) progress report \\
\path{reclassify_support_certificate.py} &
certificate reclassification utility \\
\path{exact_qls5_c12_construction.py} &
exact cardinality-\(12\) verifier \\
\path{tangent_split_c12.py} &
first-order split test around the \(c=12\) construction \\
\bottomrule
\end{tabular}
\end{center}

\end{document}
